\section{Methods}

\subsection{Data Generation}

\subsection{Cluster Training}

\subsection{Evaluation}

\subsection{Interventions}

A model trained purely on next-frame MSE could achieve deceptively strong rollout
accuracy without learning any of the underlying cellular-automaton dynamics, since
consecutive Lenia frames are highly autocorrelated and a near-identity prediction
already scores well (Section~\ref{subsection:Evaluation}). To disentangle genuine physical
understanding from simple frame-copying, we designed a set of interventions
that perturb the simulation state mid-rollout and test whether each model's
prediction still matches the behaviour of the true Lenia update rule. Concretely,
each model is first rolled out autoregressively for $t^{*}=10$ steps from a
held-out initial frame, using only its own predictions as input. The resulting
frame is then edited with one of five manually applied perturbations and fed back into
the model for a further 30 steps. The identical perturbation, with matching random
parameters (blob position, noise seed, etc.), is applied to the ground-truth Lenia
simulation at the same step, so that both trajectories start from the same
intervened state and can be compared directly at every subsequent step. We evaluate
five interventions, each targeting a different aspect of the underlying physics.
First, we inject a new "organism" (a Gaussian-blob seed) as intervention to test whether newly introduced
structures are correctly integrated. A second intervention method is zeroing a circular region to test whether
local death and regrowth are modelled. Furthermore, mirroring a small patch was applied to test sensitivity
to local symmetry breaking. To test whether global mass/intensity relationships are tracked we globally scaled the frame's density. 
Lastly, we intervened by adding pixel-wise Gaussian noise to test robustness to low-level corruption. Since no model observed
perturbed states during training, all five interventions are inherently
out-of-distribution, so a correct continuation reflects generalised understanding
of the Lenia update rule rather than memorised transitions. Results are averaged
over five active (non-dying) validation trajectories, excluding trajectories in
which the organism decays before the evaluation window ends, to avoid the
trivially low error scores produced when every model simply predicts a black
frame.